% =========================================================
% Computational Linear Algebra --- The Idea of a Linear Transformation
% (Strang \S{}8.1 / Johnston \S{}1.4 / Strang LAIA \S{}2.6)
% =========================================================

\section{The Idea of a Linear Transformation}

% ---------------------------------------------------------
\subsection*{Definition}

\begin{propositionbox}{Definition (Linear Transformation)}
A function $T : \mathbb{R}^n \to \mathbb{R}^m$ is a \emph{linear transformation} if
\[
  T(c\mathbf{v} + d\mathbf{w}) = cT(\mathbf{v}) + dT(\mathbf{w})
\]
for all $\mathbf{v}, \mathbf{w} \in \mathbb{R}^n$ and all $c, d \in \mathbb{R}$.

Equivalently, $T$ must satisfy both:
\begin{enumerate}
  \item $T(\mathbf{v} + \mathbf{w}) = T(\mathbf{v}) + T(\mathbf{w})$ \quad (preserves addition), and
  \item $T(c\mathbf{v}) = cT(\mathbf{v})$ \quad (preserves scalar multiplication).
\end{enumerate}
\end{propositionbox}

\begin{remark*}[The origin is always fixed]
Every linear transformation satisfies $T(\mathbf{0}) = \mathbf{0}$.
Therefore, \emph{translation} (shifting by a nonzero vector) is not linear.
The four geometric operations that \emph{are} linear: rotation about the origin,
scaling, reflection, and projection.
\end{remark*}

% ---------------------------------------------------------
\subsection*{\texorpdfstring{Geometric Catalog in $\mathbb{R}^2$}{Geometric Catalog in R2}}

\begin{tcolorbox}[colback=gray!6, colframe=gray!40, arc=2pt,
  left=6pt, right=6pt, top=4pt, bottom=4pt,
  title={\small\textbf{Standard Planar Transformations}},
  fonttitle=\small\bfseries]
\begin{tabular}{@{}p{0.28\textwidth}p{0.66\textwidth}@{}}
\textbf{Scaling}           & $\begin{pmatrix} c_1 & 0 \\ 0 & c_2 \end{pmatrix}$: scales $x$ by $c_1$, $y$ by $c_2$. \\[8pt]
\textbf{Rotation by $\theta$} & $\begin{pmatrix} \cos\theta & -\sin\theta \\ \sin\theta & \cos\theta \end{pmatrix}$: rotates counter-clockwise. \\[8pt]
\textbf{Reflection ($y=x$)} & $\begin{pmatrix} 0 & 1 \\ 1 & 0 \end{pmatrix}$: swaps $x$ and $y$ coordinates. \\[8pt]
\textbf{Projection ($x$-axis)} & $\begin{pmatrix} 1 & 0 \\ 0 & 0 \end{pmatrix}$: drops the $y$ component. \\
\end{tabular}
\end{tcolorbox}

\begin{remark*}[Composition is matrix multiplication]
Applying $T_1$ then $T_2$ has matrix $[T_2][T_1]$.
A sequence of geometric operations on a model (scale, then rotate, then translate)
corresponds to multiplying their matrices together in that order.
\end{remark*}

% ---------------------------------------------------------
\subsection*{Affine Transformations}

\begin{propositionbox}{Definition (Affine Transformation)}
An \emph{affine transformation} $\Phi : \mathbb{R}^n \to \mathbb{R}^n$ has the form
\[
  \Phi(\mathbf{r}) = A\mathbf{r} + \mathbf{v},
\]
where $A \in \mathbb{R}^{n \times n}$ is the linear part and $\mathbf{v} \in \mathbb{R}^n$
is a translation vector.
An affine transformation is a linear transformation when $\mathbf{v} = \mathbf{0}$.
\end{propositionbox}

\begin{remark*}[Affine invariance of B-spline curves --- NURBS Property P3.4]
An affine transformation $\Phi(\mathbf{r}) = A\mathbf{r} + \mathbf{v}$ applied to
a B-spline curve $C(u) = \sum_{i=0}^n N_{i,p}(u) P_i$ gives
\[
  \Phi(C(u)) = \sum_{i=0}^n N_{i,p}(u)\,\Phi(P_i).
\]
The transformation can be applied to the curve by applying it to each
control point individually.
This works because the basis functions sum to 1 (partition of unity):
\[
  \Phi(C(u)) = A\!\left(\sum N_{i,p} P_i\right) + \mathbf{v}
             = \sum N_{i,p} (AP_i) + \underbrace{\left(\sum N_{i,p}\right)}_{=1}\mathbf{v}
             = \sum N_{i,p}(AP_i + \mathbf{v}).
\]
This means rotating, translating, or scaling a B-spline curve costs only
$n+1$ affine operations on control points, regardless of how many points
are evaluated on the curve.
\end{remark*}
